\section{Connecting the Model to TEPL}
\label{sec:tepl-integration}
Three steps connect the theorem to the language: recognize operation
contracts, represent the mathematical construction with symbolic metadata
expressions, and verify the resulting program under the user's condition.
These steps have different proof obligations.

\subsection{From Operation Declarations to the Model}
\label{sec:declaration-adapters}
A TEPL declaration supplies an attribute schema, an \code{axes} block for
connections and output order, and a \code{semantics} block for element
values. An operation qualifies for the formal model by the structure of
this declaration, not by a name such as \code{dot} or \code{transpose}.

For contraction, check that its four fields, bindings, output order, and
value expression match Table~\ref{tab:formal-operations}. Check for extra
restrictions as well. A declaration with additional constraints is not an
unrestricted instance of the model.

Let $\rho_v$ map a declaration's fields to the model's standard roles. In
an admitted declaration, the fields are distinct and cover all roles.
Consequently the combined RHS adapter is a bijection between declared
attribute assignments and model assignments. The LHS adapter reads captured
values and preserves their sharing.

The preservation argument follows the expression tree. At an input, the
shape and values are unchanged. At an operation, assume preservation for
the children. Matching the operation contract preserves attribute legality,
the ordered output shape, and the element calculation. Induction therefore
preserves these properties for the entire expression. The bijective adapter
neither loses nor introduces attribute solutions. Recognizing the supported
templates is a finite structural check; it does not require deciding whether
an arbitrary declaration is semantically equivalent to a template.

\subsection{From Logical Axes to Symbolic Programs}
\label{sec:symbolic-bridge}
The proof uses individual port classes, finite axis sets, and sequences.
TEPL uses symbolic ordered groups and position expressions so that one rule
can describe many ranks. A symbolic interface $\mathsf I_v$ denotes an axis
order; at a concrete match it corresponds to $\operatorname{Ord}_v$.
An input interface corresponds to $I_X$. For a canonical construction,
its class set must be $P_v$, and its order must contain $Req_v$.

A mapping record $(\mathsf I,G,q)$ represents
$\operatorname{gather}(\mathsf I,q)=G$. The expression $q$ computes positions
from source metadata. Refinement splits groups while preserving their
members, order, and paired positions. Position queries implement the
$\operatorname{pos}$ operation in Equation~\eqref{eq:formal-attributes}.
Resolving intermediate queries into source-only expressions produces the
executable constructor.

These correspondences identify what a full implementation proof must show:
source extraction represents $N(L)$ faithfully; symbolic membership and
ordering operations agree with the mathematical passes; emitted expressions
evaluate to the constructed positions; and the recognized operation adapters
preserve the resulting computation. The declaration argument above covers
the adapter step. Theorem~\ref{thm:complete-synthesis} alone does not prove
all the other steps or establish completeness of symbolic proof search.

The current ordered synthesis path implements the contraction/permutation
specialization. It derives ordering requirements before construction and
leaves them as obligations under the written guard. If admission or bounded
reasoning does not support this path, TEPL can try the general synthesis
procedure in Algorithm~\ref{alg:synthesis}. Both paths check completed
candidates in a fresh context. Unary summation and the full separated-addition
extension on the exact path remain implementation work.

\subsection{Verification and Condition Suggestions}
\label{sec:condition-suggestions}
Let $C(p)$ be the user's \code{where} condition and $f$ the attribute program
used by the rule. The difference from \code{build} in
Section~\ref{sec:find} is that $f$ can be written explicitly by the user or
inferred, while \code{build} is produced by the mathematical synthesis
algorithm. Given source metadata $p$, the program computes
\begin{equation}
 \beta=f(p).
 \label{eq:attribute-program}
\end{equation}
TEPL checks Equation~\eqref{eq:contract}, requiring defined attribute
computation, RHS legality, matching output shape and dtype, and equivalent
values under the arithmetic laws permitted by the selected policy.

$C(p)$ means that evaluating \code{where} successfully returns \code{true}.
Whenever this condition holds, computing the RHS attributes must succeed
without errors, including invalid lookups, out-of-range indexing, or invalid
permutation inversions. Evaluation follows the language's ordering and
short-circuiting rules. If evaluating \code{where} fails, the rule does not
apply. An undefined $f(p)$ does not satisfy the contract.

Within the theorem's assumptions and compatible dtype and arithmetic
conditions, using its constructor reduces the obligation to
\begin{equation}
 \forall p,\quad
 \operatorname{Legal}_L(p)\land C(p)\Rightarrow\operatorname{check}(p).
 \label{eq:guard-feasibility}
\end{equation}
This is a statement about all metadata satisfying the guard. Deciding
feasibility at each concrete $p$ does not itself provide a complete proof
procedure for this universal implication. Broader operations require direct
verification of the completed RHS. The theorem concerns the mathematical
algorithm; implementation soundness is a separate obligation
(Section~\ref{sec:guarantee}).

If TEPL cannot establish sufficiency, it rejects the rule as unproved.
Failure to prove a rule does not establish that it is incorrect.
With \code{--suggest-where}, it searches for an extra condition $G(p)$,
reruns synthesis, and independently checks the new constructor under
\begin{equation}
\begin{aligned}
 \forall p,\quad &\operatorname{Legal}_L(p)\land C(p)\land G(p)\\
 &\Rightarrow p\in\operatorname{dom}(f)\land\operatorname{Valid}(p,f(p)).
\end{aligned}
 \label{eq:suggested-condition}
\end{equation}
A concrete legal metadata witness also checks that the proposed domain is
nonempty. Search is bounded: it need not find a condition, and it need not
find the weakest one. A suggestion does not accept the original rule or
silently add a runtime guard. The user must add it to \code{where}, after
which ordinary checking runs again. At runtime, verified attribute programs
use the metadata supplied by each match.

\subsection{Example: Reassociation with a Suggested Condition}
\label{sec:suggested-reassociation}
Consider this rule with omitted RHS attributes and no \code{where} block.
Imports of \code{TensorLang} and the \code{use} declaration are omitted.
\begin{codeblock}[H]
\begin{lstlisting}[language=TEPL]
rule reassociate {
  (source: dot (inner: dot A B) C)
  => (dot A (dot B C))
}
\end{lstlisting}
\end{codeblock}
The unrestricted generalized-dot rule cannot be accepted merely from
algebraic associativity: axis routing and output order also matter. Save the example in the file
named by this suggestion command:
\begin{lstlisting}[language=bash]
tepl check examples/semantic/reassociate.tepl \
  --semantic-policy=algebraic --suggest-where --rule reassociate
\end{lstlisting}
A sufficient restriction used by the reassociation example is
\begin{codeblock}[H]
\begin{lstlisting}[language=TEPL]
where {
  source.attrs.lhs_batching_dimensions == inner.axes.B.positions
}
\end{lstlisting}
\end{codeblock}
Here \code{source} and \code{inner} name operation nodes. The \code{B} in
\code{inner.axes.B} names the batch group, not input tensor \code{B}.
The condition requires the outer dot's left batch positions to be exactly
the inner dot's batch output positions, in the same order. It is the same
routing restriction used in Section~\ref{sec:partitioning}, expressed
through named nodes instead of descriptor captures.

Initially $C(p)=\top$; the proposed restriction is $G(p)$. After the user
adds this condition and checks again, synthesis fills the two omitted dot
descriptors and verification checks the result. A suggestion alone leaves
the original rule unproved. This condition is a verified sufficient
restriction, not a claim to the necessary-and-sufficient feasibility
condition of Equation~\eqref{eq:feasibility}.

\subsection{Broader Operations and Their Checking Contract}
\label{sec:broader-operators}
TEPL also supports operations outside the completeness theorem, including
broadcast, slice, convolution, reverse, negate, subtract, and
\code{zero_like}. Reverse changes element order within an axis; negate
changes signs; subtract permits cancellation; and \code{zero_like}
constructs a zero tensor of the same shape.

These operations involve coordinate transformations, negative coefficients,
or cancellation that the theorem's axis representation does not capture.
TEPL constructs candidates and independently checks their legality, ordered
output shape, and values under the written guard. For example, composing
two slices requires computing starts and strides and checking that the
combined access selects the same elements within bounds, as in
Section~\ref{sec:emit-programs}.

The intended guarantee for accepted rules is the same correctness contract,
without a completeness guarantee for discovering attributes. Failure to
synthesize or verify does not establish that no valid RHS exists. This
checking discipline is not a formal soundness proof of the current compiler;
Section~\ref{sec:guarantee} records the implementation obligations.

\subsection{Verifying Completed Rewrites}
\label{sec:correctness}
\label{sec:broader-operators}
Let $C(p)$ be the rule's \code{where} condition and $f$ its attribute program,
so that $\beta=f(p)$. Note that $f$ differs from \code{build} in
Section~\ref{sec:find}: $f$ may be user-written or inferred, whereas
\code{build} is the constructor defined by the mathematical algorithm.
Both are checked under the contract in Equation~\eqref{eq:contract}.

The verifier substitutes the attribute programs and reconstructs the completed
rewrite in a fresh context. It derives coordinate relationships and values
from those programs, rather than assuming conclusions used during synthesis.
Table~\ref{tab:verification-checks} lists the requirements. Only source
legality and the declared guard may be assumed; any additional conditions
needed by construction must follow from them.

\begin{table}[!htbp]
\caption{Checks required for a completed semantic rewrite.}
\label{tab:verification-checks}
\begin{tabularx}{\linewidth}{@{}lY@{}}\toprule
Check & Required property\\\midrule
Attribute computation & Every constructor uses only source metadata, has no
holes, and is evaluable whenever the rule applies.\\
RHS legality & Constructed attributes satisfy operation contracts, including
valid positions, paired extents, reduction domains, and access bounds.\\
Value equivalence & Fully lowered expressions agree in indexed reads,
coefficients, and reduction domains under the selected arithmetic policy.\\
Output agreement & Results have identical ordered output shapes and
compatible fixed dtype.\\\bottomrule
\end{tabularx}
\end{table}

Here $C(p)$ means that evaluating \code{where} successfully returns
\code{true}. Whenever it holds, the attribute program must evaluate without
errors, respecting the language's evaluation order and short-circuiting
rules. Undefined attribute computation cannot satisfy the contract.

Under the theorem's assumptions, using its constructor reduces the
verification obligation to
\begin{equation}
 \forall p,\quad
 \operatorname{Legal}_L(p)\land C(p)\Rightarrow\operatorname{check}(p).
 \label{eq:guard-feasibility}
\end{equation}
Whenever the source is legal and the declared condition holds,
the feasibility check must succeed.
Thus, verification must establish that the declared condition implies
the feasibility check. The theorem then guarantees correctness of
the constructed replacement for every legal match satisfying that condition.

